\documentclass[12pt]{article}
\usepackage[utf8]{inputenc}
\usepackage[T1]{fontenc}
\usepackage{hyperref}
\title{Blafe Language Resources for K-12 Curriculum Development}
\author{Compiled by Claude Sonnet 5.5 (claude-sonnet-5-5) for Andr\'as Kornai}
\date{}
\begin{document}
\maketitle

\section{Introduction}
Blafe (ISO 639-3 bfh; Glottocode blaf1238) is a Papuan language of the Western Province of Papua New Guinea, also called Tonda, Mblafe or Indorodoro, with two dialects, Mblafe and R\'anmo. It is centred on Indorodoro village, Kandarisa ward, Morehead Rural LLG, and is spoken along the Bensbach River. Speaker estimates in the sources found vary widely (about 350 to 1,200; Wikipedia gives about 670). Very little educational material appears to exist; this report lists only what the searches returned.

\section{Primary Educational Resources}
No primers, textbooks or school programmes for Blafe were found in the search results. The profile below records only that the language needs Bible translation, which implies no published vernacular literacy materials were identified.
\begin{itemize}
\item Joshua Project language profile (notes Bible translation need): \url{https://joshuaproject.net/index.php/languages/bfh}
\end{itemize}

\section{Linguistic and Reference Materials}
\begin{itemize}
\item Wikipedia, Blafe language: \url{https://en.wikipedia.org/wiki/Blafe_language}
\item New Guinea World (University of Zurich), Blafe: \url{https://newguineaworld.linguistik.uzh.ch/families/morehead-river/west-morehead-river/bensbach-river/upper-bensbach-river/blafe-ranmo/blafe}
\item Numeral systems in Blafe (Max Planck EVA): \url{https://lingweb.eva.mpg.de/channumerals/Blafe.htm}
\item ASJP word list for Blafe: \url{https://asjp.clld.org/languages/BLAFE}
\end{itemize}

\section{Religious and Community Resources}
The Joshua Project page lists Blafe as needing Bible translation; no scripture portions were found.

\section{Digital and Audio Resources}
Global Recordings Network has a page for Blafe but, according to the search summary, no recordings in Blafe itself (only related West Morehead languages).
\begin{itemize}
\item Global Recordings Network: \url{https://globalrecordings.net/en/language/bfh}
\end{itemize}

\section{Curriculum Development Context}
The language is small and under-documented in the sources found. A native-language curriculum would first require orthography work, community-driven text collection and translation of basic materials; no existing teaching materials were located. Speaker numbers should be verified locally.

\end{document}